\pagestyle{geral}
\section{Introdução}
\justifying

Este relatório documenta o desenvolvimento do terceiro e último trabalho de laboratório da unidade curricular de \gls{lsd}, intitulado "Soma de Quadrados de Números Naturais". Este trabalho sintetiza todos os conceitos abordados ao longo da disciplina, consolidando o aprendizado teórico e prático sobre circuitos digitais.

O principal objetivo é projetar e implementar um circuito digital capaz de calcular a soma dos quadrados de números naturais representados por operandos de 4 bits. Para isso, foram desenvolvidas quatro versões distintas do circuito lógico, sendo implementada apenas uma delas na placa de desenvolvimento DE10-Lite. Todas as versões serão analisadas neste relatório, com a apresentação detalhada das suas vantagens e desvantagens, justificando a escolha da versão final implementada.

Adicionalmente, o relatório abrange todos os processos relacionados à criação do circuito, incluindo a planificação, simulação, testes e implementação prática. As etapas realizadas são acompanhadas por informações detalhadas, diagramas, códigos \Gls{vhdl} e documentação necessária.

Este projeto destaca-se por integrar o design de caminhos de dados e controladores, aplicando os conhecimentos adquiridos ao longo do semestre num desafio prático. Assim, contribui para o desenvolvimento de competências essenciais em lógica digital, desde a conceção inicial até a validação final dos resultados obtidos.